\documentclass[a4paper]{article}
%To  colorize table
\usepackage[table]{xcolor}
\usepackage[a4paper, total={7in, 11.0in},includefoot, includehead, headsep=24pt, headheight=4cm]{geometry}
\usepackage{layout}%No need for this package at production.
\usepackage{setspace}%For switch between line spacing.
\usepackage{graphicx}
%For drawing lines and arbitrary shapes:
\usepackage{tikz}
%To set fixed length for table: 
\usepackage{array}
\usepackage{ifthen}
%To draw squares:
\usepackage{amssymb}
%To trim strings:
\usepackage{trimspaces}
\usepackage{ifthen}
%To draw dynamic multi-column table:
\usepackage{etoolbox}
%For references:
\usepackage[
    backend=biber,
    sorting=ydnt,
    style=ieee
]{biblatex}


\addbibresource{references.bib}

%List of needed counters:
\newcounter{length}
\newcounter{itemCount}


\title{Proposal}
\author{Name}
\date{\today}

\newcommand{\blankField}{ $\ldots\ldots\ldots\ldots\ldots$}
\makeatletter
\newcommand{\trim}[1]{\trim@spaces@noexp{#1}}
\makeatother

\usepackage{fancyhdr}
\pagestyle{fancy}
\fancyhf{}
\rhead{
    \begin{tabular}{c}
         \includegraphics[width=3cm, height=3cm]{./logo-fa-IR.png}\\
    \end{tabular}
}
\chead{
    \setstretch{2}
    بسمه تعالی \\
    دانشکده مهندسی کامپیوتر \\
    \textbf{\large فرم تعریف پروژه کارشناسی ارشد}
}

\lhead{
    \ifthenelse{\isodd{\value{page}}}
        {
            \setstretch{1.5}
            \begin{tabular}{c c}
                تاریخ: & \blankField \\
                شماره: & \blankField\\
                پیوست: & \blankField
                
                % \blankField & تاریخ: \\
                % \blankField & شماره:\\
                % \blankField & پیوست:
            \end{tabular}
        }
        {
            \setstretch{1.5}
            \begin{tabular}{c c}
                تاریخ: & \blankField \\
                شماره: & \blankField\\
                پیوست: & \blankField
                
                % \blankField & تاریخ: \\
                % \blankField & شماره:\\
                % \blankField & پیوست:
            \end{tabular}
        }
}

%Improvement: take advantages of the optional parameters! Currently I don't have any time to learn and implement it.
\newcommand{\Information}[9]{
{
    \setstretch{1.5}
    \small
    \bfseries
    \begin{tabular}{r r r}
         نام و نام خانوادگی: #1
         &
         شماره دانشجویی: #2
         &
         معدل: #3
         \\
         گرایش: #4
         &
         تعداد واحدهای گذرانده: #5
         &
         استاد راهنما: #6
         \\
         استاد راهنمای همکار: #7
         &
         تعداد واحد پروژه: #8
         &
         استاد ممتحن: #9
    \end{tabular}
}
}
\newcommand{\Title}[3]{
    {
    \setstretch{1.5}
    \bfseries
    \noindent
    عنوان کامل پروژه:
    
    فارسی: #1
    
    انگلیسی:
    \lr{#2}\\
    
    \noindent
    نوع پروژه:\;
    \begin{tabular}{m{10em} m{10em}}
        کاربردی:
        \ifthenelse{#3 = 0 \or #3 = 2}{
            $\blacksquare$
        }{
            $\square$
        }
        &
        نظری:
        \ifthenelse{#3 = 1 \or #3 = 2}{
            $\blacksquare$
        }{
            $\square$
        }
    \end{tabular}
    }
}
\newcommand{\Description}[2]{
    {
        \setstretch{1.5}
        \noindent
        \textbf{شرح مختصر پروژه:}
        \textbf{
        (با تأکيد بر اهمیت موضوع، مشکلات موجود، تعریف مسئله، کاربردها، دادگان مورد استفاده (در صورت نیاز) و نحوه ارزیابی در حدود 250 کلمه)
        }
        \noindent
        
        #1
        
        \vspace{20pt}
        
        \setcounter{length}{0}
        \setcounter{itemCount}{0}
        %First obtain the length of field
        \foreach\x in#2{%
            \addtocounter{length}{1}
        }
        
        \textbf{کلمات کلیدی:}\;
        \foreach\x in#2{%
            \addtocounter{itemCount}{1}
            \trim{\x}
            \ifthenelse{\thelength = \theitemCount}{}{-}
        }
        \vspace{1em}
    }
}

\makeatletter
\newcommand{\ProgressTable}[1]{
{
    \noindent
    \textbf{مراحل انجام پروژه و زمان‌بندی آن:}
    \\
    
    \def\tableData{}
    
    \foreach \x\y\z in #1{
        \protected@xappto\tabledata{ \x & \y & \z\\ \noexpand\hline}
    }
    
    %Add space to each row in table:
    \renewcommand{\arraystretch}{1.5}
    \begin{center}
    \begin{tabular}{| c | m{0.8\textwidth} | c |}
        \hline
        \tabledata
    \end{tabular}
    %Add space to each row in table:
        
    \end{center}
    \renewcommand{\arraystretch}{1}
}
    
}
\makeatother

\newcommand{\References}{
{
    \noindent
    \nocite{*}
    \textbf{مراجع:}
    \begin{latin}
        \printbibliography[heading=none]
    \end{latin}
}
}
\newcommand{\Courses}[1]{
{
    \noindent
    \textbf{دروس مورد نیاز:}\\
    
    %Not completed yet.
    \centering
    %Add space to each row in table:
    \renewcommand{\arraystretch}{1.5}
    \begin{center}
    \begin{tabular}{| >{\centering}m{6em} | >{\centering}m{3em} | >{\centering}m{6em} | >{\centering}m{6em} | >{\centering}m{3em} | >{\centering\arraybackslash}m{6em} |}
         \hline
         \rowcolor{lightgray}
         \multicolumn{3}{|c|}{جبرانی}
         &
         \multicolumn{3}{|c|}{تخصصی
            \tiny
            (ارتباط موضوع پروژه با دروسی که دانشجو گذرانده یا باید بگذراند)
         }
         \\
         \hline
         \rowcolor{lightgray}
         گذرانده
         &
         نمره
         &
         باید بگذراند
         &
         گذرانده
         &
         نمره
         &
         باید بگذراند
         \\
         \hline
         &
         &
        %  جبر خطی
         &
         یادگیری ماشین
         &
         20
         &
         \\
         \hline
         &
         &
         &
         تحلیل هوشمند تصاویر زیست پزشکی
         &
         20
         &
         \\
         \hline
         &
         &
         &
         &
         &
         \\
         \hline
    \end{tabular}
    \end{center}
    %Add space to each row in table:
    \renewcommand{\arraystretch}{1}
}
}
\newcommand{\Signatures}{
{
    %Add space to each row in table:
    \renewcommand{\arraystretch}{1.5}
    
    \begin{center}
        \scriptsize
        \begin{tabular}{|m{15em}|m{15em}|m{15em}|}
            \hline
            استاد راهنما:
            &
            نظر گروه:
            &
            نظر کمیته تحصیلات تکمیلی دانشکده:
            \\
            تاریخ تحویل فرم به مدیر گروه:
            &
            &
            \\
            امضای استاد راهنما:
            &
            تاریخ جلسه گروه:
            &
            تاریخ جلسه کمیته:
            \\ 
            &
            امضای مدیر گروه:
            &
            امضای معاون تحصیلات تکمیلی:
            \\
            &
            &
            \\
            &
            &
            \\
            & &\\
            & &\\
            \hline
        \end{tabular}
    \end{center}
    \footnotesize
    توجه: فرم تعریف پروژه بایستی یک روز قبل از جلسه گروه توسط استاد راهنما تحویل مدیر گروه شود.
    %Add space to each row in table:
    \renewcommand{\arraystretch}{1}
}
}

%\renewcommand{\headrulewidth}{0pt}
%\renewcommand{\footrulewidth}{0pt}

\usepackage{xepersian}
\settextfont[
Scale=1.35,
Extension=.ttf, 
Path=./,
BoldFont=bmitrabd
]{bmitra}

\begin{document}

%First parameter: Name of student(Your name)
%Second parameter: Student Number
%Third parameter: Student's GPA
%Fourth parameter: Student's Major
%Fifth parameter: Student's passed courses credits.
%Sixth parameter: Project Supervisor.
%Seventh parameter: Project Co-Supervisor
%Eighth parameter: Project Credits.
%Ninth parameter: Project Examiner
\Information{
    فرزان رحمانی
}{
    403210725
}{
    20/00
}{
    هوش مصنوعی و رباتیکز
}{
    9
}{
    دكتر محمد حسین رهبان
}{
    -
}{
    6
}{
    نام استاد ممتحن
}

\noindent

%First parameter: Title of your project in Persian.
%Second parameter: Title of your project in English.
%Third parameter: Type of your project =>
%   0: Practical.
%   1: Theoretical.
%   2: Both practical & theoretical
\Title{
    پیش‌بینی بیش‌برازش فاجعه‌بار در آموزش تخاصمی سریع
}{
    Predicting Catastrophic Overfitting in Fast Adversarial Training
}{
    2 % but maybe 2 is more suitable
}

\noindent

\Description{
‫    شبکه‌های عصبی عمیق همچنان در برابر حملات خصمانه (\lr{Adversarial Attacks}) بسیار آسیب‌پذیر هستند - اختلالات کوچک و با دقت طراحی شده که عملکرد مدل را به طور قابل توجهی کاهش می‌دهند. آموزش خصمانه (\lr{AT})، به ویژه با استفاده از روش‌های چند مرحله‌ای مانند \lr{PGD}، مؤثر اما از نظر محاسباتی پرهزینه است. آموزش خصمانه سریع (\lr{FAT})، با استفاده از \lr{FGSM} تک مرحله‌ای، سرعت و هزینه کم محاسباتی را ارائه می‌دهد اما از بیش‌برازش فاجعه‌بار (\lr{Catastrophic Overfitting}) رنج می‌برد: مدل‌ها پس از چند دوره (\lr{epoch}) آموزشی، به طور ناگهانی در برابر حملات چند مرحله‌ای، استحکام خود را از دست می‌دهند.

این پروژه پیش‌بینی کارآمد وقوع بیش‌برازش فاجعه‌بار در \lr{FAT} را بررسی می‌کند. با وجود روش‌های جلوگیری کننده موجود - مانند \lr{GradAlign}، \lr{FGSM} با نویز اضافه، \lr{N-FGSM} و \lr{AAER} - اثربخشی آنها در مجموعه‌داده‌ها و معماری‌های گوناگون متناقض است و اغلب به روش‌های اکتشافی گران‌قیمت نیاز دارند. چالش اصلی پیش‌بینی زمان و دلیل وقوع بیش‌برازش فاجعه‌بار است، نه صرفاً واکنش پس از وقوع آن.

ما پیشنهاد می‌کنیم در طول آموزش، معیارهای پیش‌بینی‌کننده‌ای ایجاد کنیم که بیش‌برازش فاجعه‌بار قریب‌الوقوع را نشان دهند. این موارد ممکن است شامل تناظر گرادیان (\lr{gradient alignment})، الگوهای گرادیان، تجزیه و تحلیل دینامیک‌های آموزش یا معیارهای انحنای مشتق شده از تجزیه و تحلیل چشم‌انداز تابع زیان (\lr{loss landscape}) باشد. با استفاده از مجموعه داده‌های معیار مانند \lr{CIFAR-10/100} و \lr{Tiny‑ImageNet} با معماری‌های استاندارد (مثلاً انواع \lr{ResNet})، این شاخص‌ها را در کنار معیارهای دقت مستحکم (\lr{robust accuracy}) پیگیری خواهیم کرد.

کاربردها: با پیش‌بینی بیش‌برازش فاجعه‌بار، الگوریتم‌های آموزشی می‌توانند به صورت پویا ابرپارامترها، میزان منظم‌سازی (\lr{regularization}) یا مراحل اختلال (\lr{perturbation stages}) را تنظیم کنند - که منجر به مدل‌های مستحکم‌تر و کارآمدتر می‌شود. این امر پیامدهای مستقیمی در کاربردهای حساس به امنیت مانند سیستم‌های خودمختار، تشخیص پزشکی و امور مالی دارد.

ارزیابی: این پروژه پیش‌بینی‌های ما از بیش‌برازش فاجعه‌بار را در برابر رویدادهای واقعی بیش‌برازش فاجعه‌بار در مجموعه‌داده‌های اعتبارسنجی، با استفاده از معیارهایی مانند میزان دقت و مقادیر تابع ضرر در حملات تک مرحله‌ای مانند \lr{FGSM} و حملات چند مرحله‌ای مانند \lr{PGD}، محک خواهد زد. موفقیت با زمان هشدار زود هنگام، کاهش وقوع بیش‌برازش فاجعه‌بار و بهبود در بده‌بستان دقت مستحکم-تمیز (\lr{robust‑clean accuracy trade-off}) اندازه‌گیری می‌شود.
‫
}{
    {
        \lr{Catastrophic Overfitting}
        ,
        \lr{Fast Adversarial Training (FAT)}
        ,
        \lr{Adversarial Robustness}
    }
}

% \hfill \break % new line
\newpage 

%Table of progress:
%   Pass line number, works and their duration in a comma delimited list.
%   line number / work / duration ,
%   example: 1 / read X / 1 year ,

\ProgressTable{
    {
    ۱/مطالعه پژوهش‌های پیشین/۳ ماه,
    ۲/بررسي نقاط ضعف روش‌ها و مدل‌های پیشین/۲ ماه,
    ۳/طراحی ايده‌های جديد و پیاده‌سازی آن‌ها/۵ ماه,
    ۴/بررسی و ارزيابی روش‌های ارائه شده/۲ ماه,
    ۵/نگارش پايان‌نامه و دفاع از پروژه/۳ ماه
    }
}

\noindent

% Add our references in references,bib
\References{}

\noindent

% Change \newcommand{\Courses} in the previous lines
\Courses{}

\noindent

\Signatures{}

\end{document}
